\documentclass[11pt]{article}

\usepackage[a4paper,margin=1in]{geometry}
\usepackage{amsmath,amssymb,mathtools}
\usepackage{microtype}
\usepackage[hidelinks]{hyperref}

\newcommand{\R}{\mathbb{R}}
\newcommand{\N}{\mathbb{N}}
\newcommand{\eps}{\varepsilon}
\newcommand{\etaapp}{\eta}
\newcommand{\Sym}{\operatorname{Sym}}
\newcommand{\diag}{\operatorname{diag}}

\title{Recurrence Formulas for Polynomial-Zonotope\\
Two-Jet Enclosures of Neural Networks}
\author{}
\date{}

\begin{document}
\maketitle

\section{Network, input, and enclosure convention}

Consider the neural network
\begin{equation}
    \Phi
    =A_L\circ\sigma_L\circ\cdots\circ A_1\circ\sigma_1\circ A_0
    :\R^{d_0}\longrightarrow\R^{d_{L+1}},
    \label{eq:network}
\end{equation}
where
\[
    A_\ell(y)=W_\ell y+b_\ell,
    \qquad
    W_\ell\in\R^{d_{\ell+1}\times d_\ell},
    \qquad \ell=0,\ldots,L,
\]
and each activation is applied componentwise,
\[
    \sigma_\ell(z)_i=\sigma(z_i),
    \qquad z\in\R^{d_\ell}.
\]

The input set is given parametrically by a polynomial zonotope
\begin{equation}
    X:[-1,1]^p\longrightarrow\R^{d_0},
    \qquad
    X(\alpha)=x_0+\sum_{\lambda\in\Lambda_X}x_\lambda\alpha^\lambda,
    \label{eq:input-pz}
\end{equation}
where \(\Lambda_X\subset\N_0^p\setminus\{0\}\) is finite and
\(\alpha^\lambda=\prod_{r=1}^p\alpha_r^{\lambda_r}\).
The variables \(\alpha\) are the \emph{domain noise symbols}.

For every activation layer \(\ell\), neuron \(i\), and derivative order
\(r\in\{0,1,2\}\), first compute a certified preactivation interval
\(I_{\ell i}\) and choose a polynomial \(\widehat\sigma_{\ell i}^{(r)}\)
and an approximation-error radius \(\rho_{\ell i}^{(r)}\geq0\) such that
\begin{equation}
    \sigma^{(r)}(t)
    \in
    \widehat\sigma_{\ell i}^{(r)}(t)
    +\rho_{\ell i}^{(r)}[-1,1]
    \qquad\text{for every }t\in I_{\ell i}.
    \label{eq:scalar-enclosure}
\end{equation}
Equivalently, at each point one may write
\begin{equation}
    \sigma^{(r)}(t)
    =\widehat\sigma_{\ell i}^{(r)}(t)
     +\rho_{\ell i}^{(r)}\eta_{\ell i}^{(r)},
    \qquad \eta_{\ell i}^{(r)}\in[-1,1].
    \label{eq:scalar-noise}
\end{equation}
These \(\eta_{\ell i}^{(r)}\) are the \emph{approximation noise symbols}.
Their coefficients are certified uniform bounds for the corresponding
residual functions.  There are no further kinds of noise symbols.

If the activation satisfies an autonomous polynomial ODE
\(\sigma'=g(\sigma)\), then
\[
    \sigma''=g'(\sigma)g(\sigma).
\]
These identities may be used to construct or certify the polynomials for the
first and second derivatives.  The recurrence below does not depend on how the
three certified scalar enclosures are obtained.

The goal is to propagate the input parametrization immediately and obtain
polynomial enclosures
\[
    Y(\eps),\qquad
    J(\eps),\qquad
    H(\eps),
    \qquad \eps=(\alpha,\eta),
\]
for
\[
    \Phi(X(\alpha)),\qquad
    \left.\partial\Phi\right|_{X(\alpha)},\qquad
    \left.\partial^2\Phi\right|_{X(\alpha)},
\]
respectively.  Thus \(J\) and \(H\) are derivatives with respect to the
physical network input \(x\in\R^{d_0}\), not derivatives with respect to
the parameter \(\alpha\).  The parameter \(\alpha\) merely records where
the two-jet is evaluated.

\section{Initialization}

Initialize the two-jet of the identity map, already evaluated on the input
polynomial zonotope, by
\begin{equation}
    Y^{\mathrm{in}}(\alpha)=X(\alpha),
    \qquad
    J^{\mathrm{in}}=I_{d_0},
    \qquad
    H^{\mathrm{in}}=0.
    \label{eq:initialization}
\end{equation}
Here
\[
    Y^{\mathrm{in}}\in\R^{d_0},\qquad
    J^{\mathrm{in}}\in\R^{d_0\times d_0},\qquad
    H^{\mathrm{in}}\in\R^{d_0\times d_0\times d_0}.
\]
All entries are polynomials in \(\alpha\); initially no approximation noise
symbols are present.

\section{The two recurrence steps}

The complete algorithm consists only of an affine step and a componentwise
activation step.

\subsection{Affine step}

Suppose \((Y,J,H)\) encloses the two-jet of a map with values in \(\R^m\),
and let \(A(y)=Wy+b\) with \(W\in\R^{n\times m}\).  Then set
\begin{equation}
    \boxed{
    \begin{aligned}
        Z_i      &= b_i+\sum_{j=1}^m W_{ij}Y_j,\\
        J^Z_{ia} &= \sum_{j=1}^m W_{ij}J_{ja},\\
        H^Z_{iab}&= \sum_{j=1}^m W_{ij}H_{jab}.
    \end{aligned}}
    \label{eq:affine-recurrence}
\end{equation}
Equivalently, \(Z=WY+b\), \(J^Z=WJ\), and \(H^Z=WH\), where \(W\)
contracts the output index of \(H\).  An affine layer introduces no new
noise symbols and no approximation error.

\subsection{Componentwise activation step}

Assume that \((Z,J^Z,H^Z)\) is the preactivation two-jet at activation layer
\(\ell\), with \(Z\in\R^{d_\ell}\).  Introduce the scalar polynomial
enclosures
\begin{equation}
    S_{\ell i}^{(r)}
    :=\widehat\sigma_{\ell i}^{(r)}(Z_i)
      +\rho_{\ell i}^{(r)}\eta_{\ell i}^{(r)},
    \qquad r=0,1,2.
    \label{eq:substituted-enclosure}
\end{equation}
Because \(Z_i\) is already a polynomial in the domain and previously
introduced approximation symbols, the substitution
\(\widehat\sigma_{\ell i}^{(r)}(Z_i)\) is again a polynomial in those symbols.

The activated two-jet is
\begin{equation}
    \boxed{
    \begin{aligned}
        Y_i
        &=S_{\ell i}^{(0)},\\
        J_{ia}
        &=S_{\ell i}^{(1)}J^Z_{ia},\\
        H_{iab}
        &=S_{\ell i}^{(2)}J^Z_{ia}J^Z_{ib}
          +S_{\ell i}^{(1)}H^Z_{iab}.
    \end{aligned}}
    \label{eq:activation-recurrence}
\end{equation}
The indices range over
\[
    i=1,\ldots,d_\ell,
    \qquad
    a,b=1,\ldots,d_0.
\]
Formula \eqref{eq:activation-recurrence} is exactly the componentwise
second-order chain rule with \(\sigma\), \(\sigma'\), and \(\sigma''\)
replaced by their certified polynomial enclosures.

The same symbol \(\eta_{\ell i}^{(1)}\), and hence the same polynomial
\(S_{\ell i}^{(1)}\), should be reused in the Jacobian formula and in the
second term of the Hessian formula.  This preserves the dependency created
by the repeated occurrence of the same quantity \(\sigma'(Z_i)\).
The symbols for derivative orders \(r=0,1,2\) remain distinct because they
certify three different scalar residuals.

For reference, in compact tensor notation let
\[
    D_\ell=\diag(S_\ell^{(1)}),
    \qquad
    E_\ell=\diag(S_\ell^{(2)}).
\]
Then
\begin{equation}
    J=D_\ell J^Z,
    \qquad
    H_{iab}
    =(E_\ell)_{ii}J^Z_{ia}J^Z_{ib}
      +(D_\ell)_{ii}H^Z_{iab}.
    \label{eq:compact-activation}
\end{equation}

\section{Layerwise algorithm}

Starting from \eqref{eq:initialization}, perform the following operations:
\begin{enumerate}
    \item Apply the affine recurrence \eqref{eq:affine-recurrence} with
          \(A_0\).
    \item For \(\ell=1,\ldots,L\):
          \begin{enumerate}
              \item apply the activation recurrence
                    \eqref{eq:activation-recurrence} for \(\sigma_\ell\);
              \item apply the affine recurrence
                    \eqref{eq:affine-recurrence} with \(A_\ell\).
          \end{enumerate}
\end{enumerate}
After the final affine map \(A_L\), the resulting arrays have shapes
\begin{equation}
    Y\in\R^{d_{L+1}},
    \qquad
    J\in\R^{d_{L+1}\times d_0},
    \qquad
    H\in\R^{d_{L+1}\times d_0\times d_0},
    \label{eq:output-shapes}
\end{equation}
and every entry is a polynomial in
\begin{equation}
    \eps=(\alpha,\eta),
    \qquad
    \alpha\in[-1,1]^p,
    \quad
    \eta\in[-1,1]^q,
    \label{eq:all-noises}
\end{equation}
for the total number \(q\) of introduced approximation symbols.  Consequently,
\((Y,J,H)\) is itself a polynomial-zonotope enclosure of the two-jet of
\(\Phi\) over \(X([-1,1]^p)\).

\section{Ground-truth identities and implementation invariants}

The recurrence is justified by the exact identities
\begin{align}
    \left.\partial(A\circ f)\right|_x
        &=W\left.\partial f\right|_x,\\
    \left.\partial^2(A\circ f)\right|_x
        &=W\left.\partial^2 f\right|_x,\\
    \left.\partial_a(\sigma\circ f)_i\right|_x
        &=\sigma'(f_i(x))\left.\partial_a f_i\right|_x,\\
    \left.\partial_{ab}(\sigma\circ f)_i\right|_x
        &=\sigma''(f_i(x))
          \left.\partial_a f_i\right|_x
          \left.\partial_b f_i\right|_x
          +\sigma'(f_i(x))\left.\partial_{ab} f_i\right|_x.
    \label{eq:exact-identities}
\end{align}
Replacing the three scalar activation quantities in these identities by
\eqref{eq:scalar-noise} gives \eqref{eq:activation-recurrence}.  Induction over
the layers therefore yields a pointwise enclosure of the value, Jacobian, and
Hessian, provided every interval \(I_{\ell i}\) contains the corresponding
preactivation range and every scalar bound \eqref{eq:scalar-enclosure} is
certified (with outward rounding in floating-point arithmetic).

The following are useful implementation invariants:
\begin{itemize}
    \item \(H_{iab}=H_{iba}\); only one triangular half need be stored if the
          remaining code respects this convention.
    \item Domain symbols are created only by the input parametrization.
          Activation layers create only approximation symbols.
    \item Polynomial coefficients are vector-, matrix-, or tensor-valued, but
          the exponent keys belong to the single common variable vector
          \(\eps=(\alpha,\eta)\).
    \item Canonicalize equal exponent vectors after polynomial products and
          substitutions; this does not alter the enclosure.
    \item Degree or term truncation is not part of the recurrence above.  If
          used, every discarded contribution must be transferred to a fresh
          certified approximation enclosure.
\end{itemize}

\end{document}
